\section{Related work}
\label{sec:related}

Closed-loop labelling sits between five literatures. The distinction this paper
turns on is the one between shortcut learning and label noise.

\paragraph{Shortcut learning.} Documented cases include
hospital-specific tokens in chest radiographs~\cite{zech2018variable}, background
rather than animal in wildlife recognition~\cite{beery2018terra}, and the general
account of shortcut learning as a property of gradient descent rather than of any
one dataset~\cite{geirhos2020shortcut, lapuschkin2019unmasking}. In each case the cue was identified after a model had already been trained and seen
to generalise badly. Crop alignment is a forward test with a value known under the
null, so it can be run before anyone has been surprised, and it names the cue in
advance.

\paragraph{Annotation artefacts.} The closest analogue is not in vision.
Hypothesis-only baselines on natural language inference showed that labels were
partly predictable from properties of the annotation process, and that leaderboard
progress had been measuring that predictability~\cite{gururangan2018annotation,
poliak2018hypothesis}. The failure is the same. The difference is that a
crowdworker's habits cannot be recovered, whereas an algorithmic labeller can be. We
invert one labeller to its constant in 17 of 17 folds and recover another labeller's
threshold and its preprocessing kernel to within 0.25\,dB on 10 of 10
events. Exact recovery lets us build the artefact with a dial instead of only detecting
it.

\paragraph{Label noise, and why this is not that.} The standard frame for imperfect
labels is noise, with a large literature on robust losses, noise-transition
estimation, and the observation that many test labels in standard benchmarks are
simply wrong~\cite{song2022learning, northcutt2021confident,
northcutt2021pervasive}. Closed-loop labels are the opposite of noise. Noise is
unpredictable from the input, so a model cannot learn it. A crop-dependent label is
deterministic in the model's own input, so a model learns it exactly. That is why
closed-loop labels inflate measured performance where noise depresses it. Treating closed-loop labels as noise
prescribes the wrong remedy, because robustness to them makes a model worse on the
benchmark and better on the truth.

\paragraph{Weak and programmatic supervision.} Labelling functions and
foundation-model pseudo-labels are the practice that produces closed-loop labels at
scale~\cite{ratner2017snorkel}. We measure one failure mode of that practice, and
the failure is about where the labelling function reads from. A rule applied once per scene is
fine. The same rule applied per crop is not.

\paragraph{Benchmark statistics.} A separate line argues that benchmark comparisons
are routinely reported without the variance or the power to support them, that seed
and data-order variation can exceed the differences being
claimed~\cite{bouthillier2021accounting, reimers2017reporting}, and that many
published effects are underpowered~\cite{card2020little, damour2022underspecification}.
Our required-sample-size framing in \Cref{sec:resolution} applies that argument to
a benchmark whose independent unit is the acquisition and not the patch, and follows
the standard practice of fixing a margin in advance~\cite{lakens2017equivalence}.

\paragraph{Spatial validation.} In remote sensing and ecology, splitting
spatially autocorrelated data at random inflates measured
skill~\cite{roberts2017cross, ploton2020spatial}, and treating correlated
observations as independent replicates has a name and a long
history~\cite{hurlbert1984pseudoreplication}. Our protocol repairs are the standard
remedies applied to a benchmark that used none of them. What we add is not the
remedies but a measurement of what skipping them costs, stated as a required sample
size.
